\documentclass[12pt]{article}

% Packages for formatting
\usepackage[margin=1in]{geometry} % Standard 1-inch margins
\usepackage{times}                % Times New Roman font (Standard for academic SOPs)
\usepackage[utf8]{inputenc}
\usepackage{setspace}             % For line spacing
\usepackage{parskip}              % Adds space between paragraphs, removes indentation

% Set line spacing to 1.15 for readability without wasting space
\setstretch{1.0}

\begin{document}

% Header Section
\begin{center}
    {\Large \textbf{Personal Statement}} \\
    \vspace{1em}
    \textbf{Applicant:} Yang Sen (Liam) Lin
\end{center}
\vspace{1em}

The defining cadence of my undergraduate years was set not in a laboratory, but in a pool. As a varsity swimmer and City-Level Gold Medalist, I committed to 8+ hours of high-intensity training weekly. This physical discipline taught me that ``oxygen debt'' is temporary, but the grit required to push through it is permanent. I applied this same resilience to a significant academic barrier: pivoting from Mechanical Engineering (ME) to Computer Science (CS).

Navigating this transition was isolating. I lacked the foundational syntax of my CS peers and often felt like an outsider in a ``gatekept'' community. However, overcoming this steep learning curve gave me a unique superpower: empathy for the uninitiated. I realized that my value wasn't just in mastering algorithms, but in dismantling the barriers that keep others out of STEM. This experience shaped my core academic identity: I am not just an engineer; I am a ``Technical Bridge.''

This philosophy transformed my approach to leadership as the Chapter Lead of the Google Developer Student Club (GDSC). I noticed that non-CS majors often felt intimidated by the rigid culture of technical academia. To counter this, I acted as a ``Translator.'' I recall mentoring a Literature major paralyzed by API logic; rather than simply correcting her syntax, I framed the code as a narrative structure with a distinct grammar. Her breakthrough reinforced my belief that diverse metaphors make technical concepts more robust. By orchestrating 10+ workshops for over 40 diverse students, I fostered an ecosystem where electrical engineers could fearlessly debate architecture with software theorists, proving that cross-disciplinary friction is fuel for innovation.

My commitment to human-centric engineering extends beyond campus walls. Initiating a meal distribution program for the homeless in Taipei was an exercise in ``human-centered logistics.'' I applied my optimization background to streamline procurement, yet quickly realized that efficiency metrics fail when they ignore human dignity. Managing these logistics grounded my high-level control theory research in gritty reality, reminding me that a system is only as robust as its ability to serve the most vulnerable.

At UCLA, I aim to continue serving as this ``Technical Bridge.'' I bring the endurance of a varsity swimmer who balances rigor with resilience, and the perspective of a leader dedicated to inclusive innovation. I plan to organize interdisciplinary roundtables---specifically connecting hardware-focused engineers with software theorists---to foster a collaborative environment where admitting ``I don't know'' is the first step toward discovery. By sharing my lived experience of navigating academic barriers, I hope to contribute to the Graduate Opportunity Program and ensure that the UCLA engineering community remains not just intellectually elite, but profoundly inclusive.

\end{document}
